\documentclass[12pt]{article}

% =========================
% Geometry & core packages
% =========================
\usepackage[a4paper,margin=0.8in]{geometry}
\usepackage{amsmath,amssymb,amsthm,mathtools}
\usepackage{bm}
\usepackage{array,booktabs}
\usepackage{enumitem}
\usepackage{microtype}
\usepackage{xcolor}
\usepackage{hyperref}
\usepackage{fancyhdr}
\usepackage{draftwatermark}
\usepackage{etoolbox}
\usepackage{titlesec}

% =========================
% Watermark
% =========================
\SetWatermarkText{Unofficial --- Refined Practice --- MH1201}
\SetWatermarkScale{1}
\SetWatermarkLightness{0.99}

% =========================
% Course metadata
% =========================
\newcommand{\CourseCode}{MH1201}
\newcommand{\CourseName}{Linear Algebra II}
\newcommand{\DocType}{Refined Practice Paper (With Solutions)}
\newcommand{\ExamMeta}{AY2025-26, Semester 2}
\newcommand{\ExamMetaShort}{Refined Practice}

% =========================
% Header/footer styles
% =========================
\fancypagestyle{main}{
\fancyhf{}
\fancyhead[L]{\CourseCode\ \CourseName\ --\ \ExamMetaShort}
\fancyhead[R]{Unofficial Practice}
\fancyfoot[C]{\thepage}
\renewcommand{\headrulewidth}{0.4pt}
\renewcommand{\footrulewidth}{0pt}
}

\fancypagestyle{plainfirst}{
\fancyhf{}
\renewcommand{\headrulewidth}{0pt}
\renewcommand{\footrulewidth}{0pt}
\fancyfoot[C]{\scriptsize\textcolor{gray}{\textit{For academic practice only. This is an original practice paper written in a similar exam style, but with different questions and full solutions.}}}
}

% =========================
% Convenience macros
% =========================
\newcommand{\dd}{\,\mathrm{d}}
\newcommand{\ee}{\mathrm{e}}
\newcommand{\ii}{\mathrm{i}}
\newcommand{\R}{\mathbb{R}}
\newcommand{\N}{\mathbb{N}}
\newcommand{\C}{\mathbb{C}}
\newcommand{\F}{\mathbb{F}}
\DeclareMathOperator{\range}{range}
\DeclareMathOperator{\Range}{Range}
\DeclareMathOperator{\Null}{Null}
\newcommand{\Span}{\operatorname{span}}
\newcommand{\dimn}{\operatorname{dim}}
\newcommand{\Ptwo}{\mathcal{P}_2(\R)}
\newcommand{\Pfour}{\mathcal{P}_4(\R)}

% =========================
% Overview block
% =========================
\newenvironment{overview}{
\par\bigskip
\begin{center}
\rule{0.92\linewidth}{0.6pt}\\[-0.2em]
\textbf{Overview}\\[-0.2em]
\rule{0.92\linewidth}{0.6pt}
\end{center}
\vspace{-0.3em}
\begin{itemize}[leftmargin=1.6em, itemsep=0.2em, topsep=0.2em]
}{
\end{itemize}
\par\bigskip
}

\setlist[enumerate]{itemsep=0.32em, topsep=0.32em}
\setlist[itemize]{itemsep=0.22em, topsep=0.22em}

\titleformat{\subsubsection}[block]
{\normalfont\large\bfseries}
{}
{0pt}
{}

\titleformat{\paragraph}[block]
{\normalfont\normalsize\bfseries}
{}
{0pt}
{}

\titlespacing*{\paragraph}
{0pt}
{1.2ex}
{0.6ex}

\setcounter{secnumdepth}{-1}
\setcounter{tocdepth}{4}

% =========================
% Title
% =========================
\title{
\Huge \CourseCode\ \CourseName \\[0.3em]
\Large \DocType \\[0.3em]
\ExamMeta
}

\date{\today}

% =========================
% Document begins
% =========================
\begin{document}

\maketitle
\thispagestyle{plainfirst}
\pagestyle{main}

\begin{overview}
\item \textbf{Time allowed:} 1 hour.
\item \textbf{Total marks:} 60.
\item \textbf{Difficulty:} Slightly harder than a standard midterm of this style.
\item \textbf{Emphasis:} Extra weight is placed on linear transformations, matrix representations, kernels, ranges, direct-sum decompositions, and basis changes.
\item \textbf{Originality policy:} The questions below are newly written and are not verbatim copies of the referenced sample midterm or tutorial sheet.
\end{overview}

\newpage
\tableofcontents
\newpage

% ============================================================
% Question 1
% ============================================================
\section{Question 1 \hfill (15 marks)}

\subsection{(a) Subspace tests}
\subsubsection{Problem}

For each of the following, determine with justification whether the given set is a vector subspace of the ambient vector space.

\begin{enumerate}[label=(\roman*)]
\item Let
\[
V\coloneqq \Pfour,
\qquad
W_1\coloneqq \{p\in \Pfour \mid p(1)=p'(0)\}.
\]
Determine whether or not \(W_1\le V\).
\hfill (3 marks)

\item Let
\[
V\coloneqq \R^{2,2},
\qquad
W_2\coloneqq \{A\in \R^{2,2}\mid A^2=A\}.
\]
Determine whether or not \(W_2\le V\).
\hfill (3 marks)

\item Let
\[
u=\begin{bmatrix}1\\2\\0\end{bmatrix},
\qquad
v=\begin{bmatrix}0\\1\\1\end{bmatrix},
\qquad
c=\begin{bmatrix}2\\5\\1\end{bmatrix},
\]
and define
\[
W_3\coloneqq \left\{au+bv+c\in \R^3 \mid a,b\in\R\right\}.
\]
Determine whether or not \(W_3\le \R^3\).
\hfill (4 marks)
\end{enumerate}

\subsubsection{Solution}

\paragraph{(i)}
Define
\[
L:\Pfour \to \R,\qquad L(p)=p(1)-p'(0).
\]
The map \(p\mapsto p(1)\) is linear, and the map \(p\mapsto p'(0)\) is also linear. Hence \(L\) is linear.

Now
\[
W_1=\{p\in \Pfour\mid L(p)=0\}=\Null(L).
\]
Since the kernel of a linear map is a subspace, it follows that
\[
\boxed{W_1\le \Pfour.}
\]

\paragraph{(ii)}
The zero matrix belongs to \(W_2\), since \(0^2=0\). So \(W_2\neq\varnothing\).

To test whether \(W_2\) is a subspace, we check closure under scalar multiplication. Let
\[
A=\begin{bmatrix}1&0\\0&0\end{bmatrix}.
\]
Then
\[
A^2=\begin{bmatrix}1&0\\0&0\end{bmatrix}=A,
\]
so \(A\in W_2\).

But
\[
(2A)^2=4A^2=4A\neq 2A.
\]
Hence \(2A\notin W_2\). Therefore \(W_2\) is not closed under scalar multiplication, so it is not a vector subspace of \(\R^{2,2}\). Thus
\[
\boxed{W_2\not\le \R^{2,2}.}
\]

\paragraph{(iii)}
We first observe that
\[
2u+v
=
2\begin{bmatrix}1\\2\\0\end{bmatrix}
+
\begin{bmatrix}0\\1\\1\end{bmatrix}
=
\begin{bmatrix}2\\5\\1\end{bmatrix}
=c.
\]
Hence \(c\in \Span\{u,v\}\).

Therefore every element of \(W_3\) can be written as
\[
au+bv+c=au+bv+(2u+v)=(a+2)u+(b+1)v.
\]
So
\[
W_3\subseteq \Span\{u,v\}.
\]
Conversely, if \(\alpha,\beta\in\R\), then
\[
\alpha u+\beta v=(\alpha-2)u+(\beta-1)v+c\in W_3.
\]
Hence
\[
\Span\{u,v\}\subseteq W_3.
\]
Therefore
\[
W_3=\Span\{u,v\},
\]
which is a subspace of \(\R^3\). Thus
\[
\boxed{W_3\le \R^3.}
\]

\subsection{(b) Linear independence under a linear map}
\subsubsection{Problem}

Let \(V\) and \(W\) be vector spaces over the same field, let \(T\in L(V,W)\), and let \(v_1,\dots,v_n\in V\). Prove that if
\[
\{T(v_1),\dots,T(v_n)\}
\]
is linearly independent in \(W\), then
\[
\{v_1,\dots,v_n\}
\]
is linearly independent in \(V\).
\hfill (5 marks)

\subsubsection{Solution}

Assume that \(\{T(v_1),\dots,T(v_n)\}\) is linearly independent in \(W\).

To prove that \(\{v_1,\dots,v_n\}\) is linearly independent in \(V\), let \(\lambda_1,\dots,\lambda_n\) be scalars such that
\[
\lambda_1 v_1+\cdots+\lambda_n v_n=0_V.
\]
Apply the linear map \(T\) to both sides. Since \(T\) is linear,
\[
T(\lambda_1 v_1+\cdots+\lambda_n v_n)
=
\lambda_1 T(v_1)+\cdots+\lambda_n T(v_n)
=
T(0_V)=0_W.
\]
Hence
\[
\lambda_1 T(v_1)+\cdots+\lambda_n T(v_n)=0_W.
\]
But the vectors \(T(v_1),\dots,T(v_n)\) are linearly independent, so
\[
\lambda_1=\cdots=\lambda_n=0.
\]
Therefore \(v_1,\dots,v_n\) are linearly independent. Thus
\[
\boxed{\{v_1,\dots,v_n\}\text{ is linearly independent in }V.}
\]

\newpage

% ============================================================
% Question 2
% ============================================================
\section{Question 2 \hfill (18 marks)}

\subsection{Problem}

Define a function \(T:\R^{2,2}\to \R^{2,2}\) by
\[
\forall A\in \R^{2,2}:\qquad
T(A)\coloneqq A+A^{\mathsf T}-\operatorname{tr}(A)I_2.
\]

\begin{enumerate}[label=(\alph*)]
\item \textbf{[3 points]} Show that \(T\in L(\R^{2,2})\).

\item \textbf{[4 points]} Compute the standard matrix representation of \(T\) with respect to the ordered basis
\[
(E_{11},E_{12},E_{21},E_{22}).
\]

\item \textbf{[4 points]} Derive a basis for \(\Null(T)\).

\item \textbf{[4 points]} Derive a basis for \(\Range(T)\).

\item \textbf{[3 points]} Determine, with justification, whether or not
\[
\R^{2,2}=\Null(T)\oplus \Range(T).
\]
Your argument must include an explicit decomposition of an arbitrary matrix \(B\in \R^{2,2}\).
\end{enumerate}

\subsubsection{Solution}

Let
\[
A=\begin{bmatrix}a&b\\c&d\end{bmatrix}.
\]
Then
\[
A^{\mathsf T}=\begin{bmatrix}a&c\\b&d\end{bmatrix},
\qquad
\operatorname{tr}(A)=a+d,
\qquad
\operatorname{tr}(A)I_2=
\begin{bmatrix}a+d&0\\0&a+d\end{bmatrix}.
\]
Hence
\[
T(A)
=
\begin{bmatrix}a&b\\c&d\end{bmatrix}
+
\begin{bmatrix}a&c\\b&d\end{bmatrix}
-
\begin{bmatrix}a+d&0\\0&a+d\end{bmatrix}
=
\begin{bmatrix}
a-d & b+c\\
b+c & d-a
\end{bmatrix}.
\]

\begin{enumerate}[label=(\alph*)]
\item \textbf{Linearity.}

The maps
\[
A\mapsto A,\qquad A\mapsto A^{\mathsf T},\qquad A\mapsto \operatorname{tr}(A)I_2
\]
are all linear. Therefore
\[
T(A)=A+A^{\mathsf T}-\operatorname{tr}(A)I_2
\]
is a linear combination of linear maps, hence is linear. Thus
\[
\boxed{T\in L(\R^{2,2}).}
\]

\item \textbf{Standard matrix representation.}

Using the standard basis \(E_{11},E_{12},E_{21},E_{22}\), we compute:
\[
T(E_{11})
=
\begin{bmatrix}1&0\\0&-1\end{bmatrix},
\qquad
T(E_{12})
=
\begin{bmatrix}0&1\\1&0\end{bmatrix},
\]
\[
T(E_{21})
=
\begin{bmatrix}0&1\\1&0\end{bmatrix},
\qquad
T(E_{22})
=
\begin{bmatrix}-1&0\\0&1\end{bmatrix}.
\]
Identifying
\[
\begin{bmatrix}x&y\\z&w\end{bmatrix}\longleftrightarrow (x,y,z,w),
\]
the columns are
\[
[T(E_{11})]=(1,0,0,-1),\quad
[T(E_{12})]=(0,1,1,0),\quad
[T(E_{21})]=(0,1,1,0),\quad
[T(E_{22})]=(-1,0,0,1).
\]
Hence
\[
[T]=
\begin{bmatrix}
1&0&0&-1\\
0&1&1&0\\
0&1&1&0\\
-1&0&0&1
\end{bmatrix}.
\]
Therefore
\[
\boxed{
[T]=
\begin{bmatrix}
1&0&0&-1\\
0&1&1&0\\
0&1&1&0\\
-1&0&0&1
\end{bmatrix}.}
\]

\item \textbf{Basis for \(\Null(T)\).}

We solve \(T(A)=0\). From
\[
T(A)=
\begin{bmatrix}
a-d & b+c\\
b+c & d-a
\end{bmatrix},
\]
we get the system
\[
a-d=0,\qquad b+c=0.
\]
Thus
\[
d=a,\qquad c=-b.
\]
So every matrix in the null space has the form
\[
\begin{bmatrix}a&b\\-b&a\end{bmatrix}
=
a\begin{bmatrix}1&0\\0&1\end{bmatrix}
+
b\begin{bmatrix}0&1\\-1&0\end{bmatrix}.
\]
Hence
\[
\Null(T)=
\Span\left\{
\begin{bmatrix}1&0\\0&1\end{bmatrix},
\begin{bmatrix}0&1\\-1&0\end{bmatrix}
\right\}.
\]
Therefore a basis for \(\Null(T)\) is
\[
\boxed{
\left\{
\begin{bmatrix}1&0\\0&1\end{bmatrix},
\begin{bmatrix}0&1\\-1&0\end{bmatrix}
\right\}.}
\]

\item \textbf{Basis for \(\Range(T)\).}

From the formula
\[
T(A)=
\begin{bmatrix}
a-d & b+c\\
b+c & d-a
\end{bmatrix},
\]
every output has the form
\[
\begin{bmatrix}x&y\\y&-x\end{bmatrix}.
\]
So
\[
\Range(T)\subseteq
\left\{
\begin{bmatrix}x&y\\y&-x\end{bmatrix}
:x,y\in\R
\right\}.
\]
Conversely, given any matrix
\[
Y=\begin{bmatrix}x&y\\y&-x\end{bmatrix},
\]
choose
\[
A=\begin{bmatrix}x&y\\0&0\end{bmatrix}.
\]
Then
\[
T(A)=
\begin{bmatrix}
x&y\\
y&-x
\end{bmatrix}
=Y.
\]
Thus
\[
\Range(T)=
\left\{
\begin{bmatrix}x&y\\y&-x\end{bmatrix}
:x,y\in\R
\right\}.
\]
Hence
\[
\Range(T)
=
\Span\left\{
\begin{bmatrix}1&0\\0&-1\end{bmatrix},
\begin{bmatrix}0&1\\1&0\end{bmatrix}
\right\}.
\]
Therefore a basis for \(\Range(T)\) is
\[
\boxed{
\left\{
\begin{bmatrix}1&0\\0&-1\end{bmatrix},
\begin{bmatrix}0&1\\1&0\end{bmatrix}
\right\}.}
\]

\item \textbf{Direct sum decomposition.}

Let
\[
B=\begin{bmatrix}p&q\\r&s\end{bmatrix}\in \R^{2,2}.
\]
Define
\[
R=
\frac12
\begin{bmatrix}
p-s & q+r\\
q+r & s-p
\end{bmatrix},
\qquad
N=
\frac12
\begin{bmatrix}
p+s & q-r\\
r-q & p+s
\end{bmatrix}.
\]
Then
\[
B=R+N.
\]

We check that \(R\in \Range(T)\). Indeed, \(R\) has the form
\[
\begin{bmatrix}x&y\\y&-x\end{bmatrix},
\]
so \(R\in \Range(T)\).

We also check that \(N\in \Null(T)\). Indeed, \(N\) has the form
\[
\begin{bmatrix}a&b\\-b&a\end{bmatrix},
\]
so \(N\in \Null(T)\).

Hence every \(B\in \R^{2,2}\) can be written as
\[
B=R+N
\qquad\text{with }R\in \Range(T),\ N\in \Null(T).
\]
So
\[
\R^{2,2}=\Range(T)+\Null(T).
\]

Next, suppose \(X\in \Range(T)\cap \Null(T)\). Then \(X\) must have both forms
\[
X=\begin{bmatrix}x&y\\y&-x\end{bmatrix}
\qquad\text{and}\qquad
X=\begin{bmatrix}a&b\\-b&a\end{bmatrix}.
\]
Comparing entries gives
\[
x=a,\qquad y=b,\qquad y=-b,\qquad -x=a.
\]
Hence
\[
y=b=0,\qquad x=a=0.
\]
So \(X=0\). Therefore
\[
\Range(T)\cap \Null(T)=\{0\}.
\]

Thus
\[
\boxed{\R^{2,2}=\Null(T)\oplus \Range(T).}
\]
\end{enumerate}

\newpage

% ============================================================
% Question 3
% ============================================================
\section{Question 3 \hfill (15 marks)}

\subsection{Problem}

Let \(\beta\) be the standard ordered basis of \(\Ptwo\), namely
\[
\beta=(1,X,X^2),
\]
and let
\[
\gamma=(1,\ X-1,\ (X-1)^2).
\]
Define \(S:\Ptwo\to \Ptwo\) by
\[
\forall p\in \Ptwo:\qquad S(p)\coloneqq p+(X-1)p'.
\]

\begin{enumerate}[label=(\alph*)]
\item \textbf{[3 points]} Compute the change-of-coordinates matrix
\[
[I_{\Ptwo}]^\beta_\gamma.
\]

\item \textbf{[5 points]} Compute
\[
[S]^\gamma_\gamma.
\]

\item \textbf{[4 points]} Determine whether \(S\) is injective and whether \(S\) is surjective.

\item \textbf{[3 points]} Find the unique polynomial \(p\in \Ptwo\) such that
\[
S(p)=X^2+1.
\]
\end{enumerate}

\subsubsection{Solution}

\begin{enumerate}[label=(\alph*)]
\item \textbf{Compute \([I_{\Ptwo}]^\beta_\gamma\).}

The columns of \([I_{\Ptwo}]^\beta_\gamma\) are the \(\beta\)-coordinates of the vectors in \(\gamma\).

First,
\[
1=1,
\]
so
\[
[1]_\beta=\begin{bmatrix}1\\0\\0\end{bmatrix}.
\]

Next,
\[
X-1=-1+X,
\]
so
\[
[X-1]_\beta=\begin{bmatrix}-1\\1\\0\end{bmatrix}.
\]

Finally,
\[
(X-1)^2=X^2-2X+1,
\]
so
\[
[(X-1)^2]_\beta=\begin{bmatrix}1\\-2\\1\end{bmatrix}.
\]

Therefore
\[
[I_{\Ptwo}]^\beta_\gamma
=
\begin{bmatrix}
1&-1&1\\
0&1&-2\\
0&0&1
\end{bmatrix}.
\]
Hence
\[
\boxed{
[I_{\Ptwo}]^\beta_\gamma
=
\begin{bmatrix}
1&-1&1\\
0&1&-2\\
0&0&1
\end{bmatrix}.}
\]

\item \textbf{Compute \([S]^\gamma_\gamma\).}

Let
\[
u=X-1.
\]
Then the ordered basis \(\gamma\) is
\[
(1,u,u^2).
\]
We compute \(S\) on these basis vectors.

Since \(1'=0\),
\[
S(1)=1+u\cdot 0=1.
\]

Since \(u'=1\),
\[
S(u)=u+u\cdot 1=2u.
\]

Since \((u^2)'=2u\),
\[
S(u^2)=u^2+u(2u)=3u^2.
\]

Thus
\[
[S(1)]_\gamma=\begin{bmatrix}1\\0\\0\end{bmatrix},\qquad
[S(u)]_\gamma=\begin{bmatrix}0\\2\\0\end{bmatrix},\qquad
[S(u^2)]_\gamma=\begin{bmatrix}0\\0\\3\end{bmatrix}.
\]
Therefore
\[
[S]^\gamma_\gamma=
\begin{bmatrix}
1&0&0\\
0&2&0\\
0&0&3
\end{bmatrix}.
\]
Hence
\[
\boxed{
[S]^\gamma_\gamma=
\begin{bmatrix}
1&0&0\\
0&2&0\\
0&0&3
\end{bmatrix}.}
\]

\item \textbf{Determine injectivity and surjectivity.}

The matrix
\[
[S]^\gamma_\gamma=
\begin{bmatrix}
1&0&0\\
0&2&0\\
0&0&3
\end{bmatrix}
\]
has a pivot in every column and a pivot in every row. Therefore \(S\) is injective and surjective.

Equivalently, since the diagonal entries are all nonzero, \([S]^\gamma_\gamma\) is invertible, so \(S\) is invertible. Thus
\[
\boxed{S\text{ is injective and surjective.}}
\]

\item \textbf{Solve \(S(p)=X^2+1\).}

Again let \(u=X-1\). Then
\[
X=u+1,
\qquad
X^2+1=(u+1)^2+1=u^2+2u+2.
\]
So in the basis \(\gamma=(1,u,u^2)\),
\[
[X^2+1]_\gamma=\begin{bmatrix}2\\2\\1\end{bmatrix}.
\]

Suppose
\[
[p]_\gamma=\begin{bmatrix}a\\b\\c\end{bmatrix}.
\]
Then
\[
[S(p)]_\gamma=
[S]^\gamma_\gamma [p]_\gamma
=
\begin{bmatrix}
1&0&0\\
0&2&0\\
0&0&3
\end{bmatrix}
\begin{bmatrix}a\\b\\c\end{bmatrix}
=
\begin{bmatrix}a\\2b\\3c\end{bmatrix}.
\]
We want
\[
[S(p)]_\gamma=\begin{bmatrix}2\\2\\1\end{bmatrix}.
\]
Hence
\[
a=2,\qquad 2b=2,\qquad 3c=1,
\]
so
\[
a=2,\qquad b=1,\qquad c=\frac13.
\]
Therefore
\[
p=2+(X-1)+\frac13(X-1)^2.
\]
Expanding,
\[
p
=
2+X-1+\frac13(X^2-2X+1)
=
\frac13X^2+\frac13X+\frac43.
\]
Thus the unique polynomial is
\[
\boxed{
p(X)=2+(X-1)+\frac13(X-1)^2
=
\frac13X^2+\frac13X+\frac43.}
\]
\end{enumerate}

\newpage

% ============================================================
% Question 4
% ============================================================
\section{Question 4 \hfill (12 marks)}

\subsection{Problem}

Let \(V\) be a vector space over a field \(F\) with \(\operatorname{char}(F)\neq 2\), and let \(S\in L(V)\) satisfy
\[
S^2=I_V.
\]
Define
\[
U\coloneqq \{v\in V\mid S(v)=v\},
\qquad
W\coloneqq \{v\in V\mid S(v)=-v\}.
\]

\begin{enumerate}[label=(\roman*)]
\item \textbf{[3 points]} Show that \(U\) and \(W\) are vector subspaces of \(V\).

\item \textbf{[3 points]} Show that
\[
U\cap W=\{0_V\}.
\]

\item \textbf{[6 points]} Prove that
\[
V=U\oplus W.
\]
Your proof must give an explicit decomposition of an arbitrary \(v\in V\).
\end{enumerate}

\subsubsection{Solution}

\paragraph{(i) \(U\) and \(W\) are subspaces.}
Define
\[
L_+=S-I_V,\qquad L_-=S+I_V.
\]
Since \(S\) and \(I_V\) are linear, both \(L_+\) and \(L_-\) are linear.

Now
\[
U=\{v\in V\mid S(v)=v\}=\{v\in V\mid (S-I_V)(v)=0\}=\Null(L_+),
\]
and
\[
W=\{v\in V\mid S(v)=-v\}=\{v\in V\mid (S+I_V)(v)=0\}=\Null(L_-).
\]
Since kernels of linear maps are subspaces, both \(U\) and \(W\) are vector subspaces of \(V\). Therefore
\[
\boxed{U\le V\quad\text{and}\quad W\le V.}
\]

\paragraph{(ii) Trivial intersection.}
Let \(x\in U\cap W\). Then
\[
S(x)=x
\qquad\text{and}\qquad
S(x)=-x.
\]
Hence
\[
x=-x.
\]
So
\[
2x=0.
\]
Since \(\operatorname{char}(F)\neq 2\), the scalar \(2\) is nonzero in \(F\), and therefore \(x=0\). Thus
\[
U\cap W=\{0_V\}.
\]
Hence
\[
\boxed{U\cap W=\{0_V\}.}
\]

\paragraph{(iii) Direct-sum decomposition.}
Let \(v\in V\) be arbitrary. Define
\[
u\coloneqq \frac12(v+S(v)),
\qquad
w\coloneqq \frac12(v-S(v)).
\]
Then
\[
u+w
=
\frac12(v+S(v))+\frac12(v-S(v))
=
v.
\]
So it remains to show that \(u\in U\) and \(w\in W\).

First,
\[
S(u)
=
S\left(\frac12(v+S(v))\right)
=
\frac12(S(v)+S^2(v))
=
\frac12(S(v)+v)
=
u,
\]
since \(S^2=I_V\). Hence \(u\in U\).

Next,
\[
S(w)
=
S\left(\frac12(v-S(v))\right)
=
\frac12(S(v)-S^2(v))
=
\frac12(S(v)-v)
=
-\frac12(v-S(v))
=
-w.
\]
Hence \(w\in W\).

Therefore every \(v\in V\) can be written as
\[
v=u+w
\qquad\text{with }u\in U,\ w\in W.
\]
So
\[
V=U+W.
\]
Together with part (ii), which showed that \(U\cap W=\{0_V\}\), we conclude that
\[
\boxed{V=U\oplus W.}
\]

\newpage

% ============================================================
% Brief mark scheme
% ============================================================
\section{Brief Mark Scheme}

\subsection{Question 1 \hfill (15 marks)}
\begin{itemize}[leftmargin=1.6em]
\item \textbf{(a)(i) (3 marks)} Correct identification of a linear functional / kernel argument and conclusion.
\item \textbf{(a)(ii) (3 marks)} Correct counterexample to closure under scalar multiplication or addition and conclusion.
\item \textbf{(a)(iii) (4 marks)} Correct verification that \(c\in\Span\{u,v\}\), deduction that \(W_3=\Span\{u,v\}\), and conclusion.
\item \textbf{(b) (5 marks)} Correct linear-independence proof using linearity of \(T\).
\end{itemize}

\subsection{Question 2 \hfill (18 marks)}
\begin{itemize}[leftmargin=1.6em]
\item \textbf{(a) (3 marks)} Correct linearity argument.
\item \textbf{(b) (4 marks)} Correct images of basis matrices and assembly of \([T]\).
\item \textbf{(c) (4 marks)} Correct solution of \(T(A)=0\) and a valid basis of \(\Null(T)\).
\item \textbf{(d) (4 marks)} Correct characterisation and basis of \(\Range(T)\).
\item \textbf{(e) (3 marks)} Explicit decomposition, trivial intersection, and direct-sum conclusion.
\end{itemize}

\subsection{Question 3 \hfill (15 marks)}
\begin{itemize}[leftmargin=1.6em]
\item \textbf{(a) (3 marks)} Correct change-of-coordinates matrix.
\item \textbf{(b) (5 marks)} Correct computation of \([S]^\gamma_\gamma\).
\item \textbf{(c) (4 marks)} Correct injective/surjective conclusion with justification.
\item \textbf{(d) (3 marks)} Correct preimage polynomial.
\end{itemize}

\subsection{Question 4 \hfill (12 marks)}
\begin{itemize}[leftmargin=1.6em]
\item \textbf{(i) (3 marks)} Correct proof that \(U,W\) are subspaces.
\item \textbf{(ii) (3 marks)} Correct proof that \(U\cap W=\{0_V\}\).
\item \textbf{(iii) (6 marks)} Correct explicit decomposition \(v=\frac12(v+Sv)+\frac12(v-Sv)\) and direct-sum conclusion.
\end{itemize}

\end{document}
